\chapter{Descrição, cobertura e seleção textual}\label{cap:descritiva}
\section{Comparação dos universos}
A Tabela~\ref{tab:descritivas} separa catálogo, amostra e domínio aceito, com e sem pesos quando pertinente. A popularidade média da amostra ponderada é próxima da média do catálogo, enquanto as faixas aceitas apresentam média mais alta. A proximidade entre amostra e catálogo é uma propriedade descritiva do sorteio registrado; não implica que toda variável tenha sido controlada por estratificação nem que a etapa textual preserve o desenho.

\begin{table}[htbp]
\centering\small
\caption{Popularidade observada: catálogo, amostra e domínio com letras aceitas.}
\label{tab:descritivas}
\begin{tabular}{lrrrrr}
\toprule
Universo / avaliação & n & Média & P25 & Mediana & P75 \\
\midrule
Catálogo & 1.159.764 & 18,3831 & 5 & 15 & 29 \\
Amostra, sem pesos & 200.000 & 18,6752 & 5 & 15 & 29 \\
Amostra, ponderada & 200.000 & 18,3899 & 5 & 15 & 29 \\
Aceitas, sem pesos & 77.535 & 23,9821 & 12 & 22 & 34 \\
Aceitas, ponderadas & 77.535 & 23,7482 & 12 & 22 & 34 \\
\bottomrule
\end{tabular}
\fonte{Elaboração própria a partir dos agregados auditados do estudo.}
\nota{Quantis pela inversa da distribuição empírica, ponderada quando indicado. n conta faixas, não soma de pesos.}
\end{table}


A média sem pesos das aceitas é 23,9821, contra 18,3831 no catálogo. A mediana passa de 15 para 22 e o terceiro quartil de 29 para 34. A Figura~\ref{fig:distribuicao} apresenta distribuições acumuladas calculadas a partir das frequências reais. O deslocamento das aceitas para valores maiores é compatível com disponibilidade textual seletiva. Não identifica sua causa, nem prova que letras provoquem popularidade.

\figura{distribuicao_popularidade}{Distribuição acumulada da popularidade nos universos definidos.}{fig:distribuicao}{Agregados da auditoria do catálogo, amostra estratificada e parecer da recuperação v2. A curva das aceitas usa pesos iguais.}
\FloatBarrier

\section{Cobertura textual}
A aceitação v2 corresponde a 38,7675\% das 200.000 faixas. A razão com pesos originais é 38,7436\%, obtida pela soma dos pesos das aceitas dividida pela soma dos pesos da amostra inteira. Essa proximidade global entre proporções não elimina diferenças de composição dentro do domínio aceito. O maior grupo restante é o de não encontradas, seguido das ambiguidades.

A Figura~\ref{fig:cobertura} descreve a aceitação por banda de popularidade na amostra. As 499 faixas com popularidade pelo menos 80 foram incluídas no sorteio, mas apenas 128 tiveram letra aceita: inclusão amostral e observação textual são eventos distintos. Mesmo na banda integral, não se pode usar a existência da amostra para afirmar cobertura textual completa.

\figura{cobertura_bandas}{Proporção de letras aceitas por banda de popularidade.}{fig:cobertura}{Contagens da recuperação v2, com denominador da amostra em cada banda; porcentagens sem pesos.}
\FloatBarrier

O parecer de cobertura registrou 102 estratos sem nenhuma letra aceita e 348 com menos de cinco aceitações. A inexistência ou escassez de casos em células do desenho limita reponderações e comparações específicas. Um peso original positivo não cria informação textual em um estrato não observado.

\section{Seleção e interpretação das marginais}
A avaliação de compatibilidade das aceitas com o catálogo examinou 104 critérios primários ponderados, com tolerâncias operacionais explícitas. Foram 21 resultados de equivalência, 74 diferenças individuais e nove inconclusivos; 72 diferenças tiveram apoio no controle simultâneo adotado. Esses diagnósticos não autorizam chamar a base aceita de representativa do catálogo.

As margens foram de um ponto para a média de popularidade, 0,10 desvio do catálogo para outras médias contínuas e year, e o menor entre um ponto percentual e 20\% da proporção do catálogo para categorias. O parecer utilizou linearização de razões sob o desenho estratificado e TOST de Wald a 5\%, associado a limites de 90\% por comparação, além de limites de 95\% e controle Bonferroni para diferenças simultâneas. A equivalência conjunta exigia todos os critérios, sem escolher uma maioria favorável. Células raras ou de variância não estimável foram classificadas como inconclusivas. São escolhas e aproximações do diagnóstico registrado, não propriedades garantidas da base.

As tolerâncias e comparações foram formuladas depois de etapas anteriores do projeto e não constituem confirmação independente. Além disso, falta de equivalência não é sinônimo de diferença demonstrada: o parecer distingue esses resultados por critério. Tampouco equivalência marginal, se ocorresse, certificaria identidade correta das letras, equivalência multivariada ou ausência de viés preditivo.

Na análise principal, os pesos originais não foram usados para tentar corrigir a disponibilidade de letras. Eles permanecem preservados para descrição e sensibilidade. Uma correção inferencial para a seleção exigiria hipóteses sobre observação, suporte e especificação que não são verificadas apenas pelos agregados existentes.

\section{Idioma, repetição e qualidade dos preditores}
O domínio permaneceu multilingue. O detector automático classificou 56.908 faixas como inglesas, 20.447 como não inglesas e 180 como incertas. Esses números se referem ao status que combina código estimado e confiança; não devem ser confundidos com a simples contagem de um código de língua. O detector não foi tratado como verdade humana nem usado para exclusão silenciosa.

Foram registrados 484 vínculos excedentes de faixas em relação às entradas normalizadas únicas. Conteúdo compartilhado é relevante para a avaliação porque pode aproximar versões ou ligar artistas. Por isso, duplicação textual foi incorporada ao agrupamento, sem eliminar a faixa da distribuição empírica de interesse.

Não houve ausência ou invalidade nos 15 campos do Baseline na coorte final. Essa constatação não valida todos os domínios possíveis da fonte nem corrige seu mecanismo de coleta. A auditoria do catálogo tinha escopo mais amplo e regras diagnósticas próprias. Os controles de elegibilidade, qualidade e modelagem não devem ser intercambiados retrospectivamente para excluir observações difíceis.
